\subsection{Income and purchase costs}
The hypothetical road shock raises the benchmark population-weighted income
measure by 3.2 log points and the cost of initial purchases by 1.6. Their difference
gives a local welfare gain of 1.6 log points. This accounting comparison values
both terms for the same population and shock. It does not by itself show whether
the income response is recovered correctly from spending or whether every
resident gains. Comparing income restrictions is therefore part of interpreting
the result, rather than a separate exercise with no bearing on the benchmark mean.

\subsection{Sensitivity of income recovery}
An alternative allocation gives income gains of 2.4 log points with the same
initial purchase cost change, leaving welfare at 0.8. Changing the allocation changes
how receipts are assigned in the calculation; it does not establish a change in
their actual recipients. A separate unrestricted recovery is more consequential:
income falls by 0.5 log points and welfare by 2.1, again with costs rising by 1.6.
Similar spatial rankings across these exercises preserve part of the relative
comparison, but not the positive mean. They are not confidence bounds, and the
supplied evidence does not rank the credibility of the restrictions. The
comparison with observed income cannot settle the disagreement over the common income
change because month effects absorb that component. It provides a comparison
of spatial responses, not an independent estimate of aggregate income.

\subsection{Structural and policy comparisons}
The structural exercise asks how the same shock changes welfare when prices
and choices adjust jointly in the model. Solving the model with fixed workplaces
gives a finite gain of 2.6 percent, compared with 2.5 percent when workers choose
workplaces. Both use the same baseline and shock, and total hours remain fixed.
The mobile alternative changes workplace choice rather than adding an hours
response. These finite percentages are different objects from the empirical
local calculation in log points, which was not a calibration target. Similar model
means neither establish similar gains for each residence nor resolve the
disagreement over income recovery in the spending estimates.

The policy exercise yields gains of 0.30 percent under a uniform tax and 0.34
percent under differentiated taxes. The additional gain is 0.04 percentage points;
it is not 0.04 percent of the gain under the uniform policy. This modest additional gain concerns
the specified feasible policies and model. It does not imply that differences in
local rates are caused by sparse data. A policy with constant rates within districts
has not yet been computed, and checks of local necessary conditions do not establish
a unique global optimum. The policy results should therefore answer the model's
tax question without becoming a recommendation based on observed data alone.

Together these comparisons locate the stronger and weaker conclusions. Much of
the spatial ordering remains similar across changes in recovery, but the mean welfare gain is
conditional on the income restriction. Structural comparisons show what changes
under additional assumptions about adjustment and policy, but cannot replace
evidence for that empirical restriction. The economic answer is thus a spatial
incidence pattern with a conditional aggregate gain, not a single welfare number
validated by several close model predictions.
